% GENERATED FILE. Edit canonical lessons, then run npm run generate:books.
% canonical-source-hash: fnv1a64:2322b4d18e75e96e
% canonical-lessons: LA-C12-medius-dies-nox

\chapter{Noon and Midnight}
\label{ch:la-noon-and-midnight}
\hlchaptermodality{12}

\begin{chapteropening}
\textbf{By the end of this chapter:} \emph{I can say noon and midnight in Latin, and explain why English's noon does not come from either.}

Say medius diēs and media nox, then follow both halves into Spanish and French, and trace English noon to an unrelated \textquotedblleft{}ninth hour\textquotedblright{}.
\end{chapteropening}

\section[medius diēs, media nox]{medius diēs, media nox — the source phrase, worn down at different rates}

\label{lesson:LA-C12-medius-dies-nox}



\begin{quote}
Here's the same Latin phrase, followed into three languages — and each one treated it completely differently.
\end{quote}

\subsection*{You'll want to know: The two phrases}
\begin{itemize}
\raggedright
  \item \textbf{medius diēs} — literally \textquotedblleft{}\textbf{the middle day},\textquotedblright{} meaning \textquotedblleft{}\textbf{midday, noon}\textquotedblright{} (\emph{medius}, \textquotedblleft{}middle\textquotedblright{} + \emph{diēs}, \textquotedblleft{}day\textquotedblright{}).
  \item \textbf{media nox} — literally \textquotedblleft{}\textbf{the middle night},\textquotedblright{} meaning \textquotedblleft{}\textbf{midnight}\textquotedblright{} (\emph{media}, the feminine of \emph{medius} + \emph{nox, noctis}, \textquotedblleft{}night\textquotedblright{}).
\end{itemize}

\begin{culture}
\begin{itemize}
\raggedright
  \item \textbf{Spanish} kept \textbf{both halves of both phrases} as fully ordinary, living words: \emph{medio/media} (\textquotedblleft{}middle\textquotedblright{}) and \emph{día/noche} (\textquotedblleft{}day/night\textquotedblright{}) — giving \emph{mediodía} and \emph{medianoche}, entirely transparent.
  \item \textbf{French} eroded the \textbf{same phrase unevenly}. \emph{Nox} survived as the everyday word \emph{nuit} (\textquotedblleft{}night\textquotedblright{}) — still fully transparent inside \emph{minuit}. But \emph{diēs} did \textbf{not} survive as French's everyday word for \textquotedblleft{}day\textquotedblright{} (that's \emph{jour} instead) — so inside \emph{midi}, the \emph{-di} fragment is a fossil, recognizable only from weekday names like \emph{lundi, mardi}, not from any living, standalone word for \textquotedblleft{}day.\textquotedblright{}
  \item \textbf{English} didn't inherit this phrase for \textquotedblleft{}noon\textquotedblright{} \textbf{at all}. Instead, English's \emph{noon} comes from an entirely different Latin phrase, \textquotedblleft{}the \textbf{ninth hour}\textquotedblright{} (the next lesson meets its second word) — originally meaning roughly 3 p.m. (counting from a sunrise start), it drifted steadily earlier over the centuries until it settled on \textbf{midday}.
\end{itemize}

One Latin phrase, three completely different fates: fully kept, half-worn, and simply replaced by another Latin phrase altogether.
\end{culture}

\begin{cousinweb}
\textbf{medius diēs} joins two inherited words. \emph{medius} \textquotedblleft{}middle\textquotedblright{} is usually traced to Proto-Indo-European \emph{*medhyos}, \textquotedblleft{}middle\textquotedblright{} (EtymDB 2.1), a cousin of English mid. You met \emph{diēs} \textquotedblleft{}day\textquotedblright{} with the weekday names; \emph{nox} \textquotedblleft{}night\textquotedblright{} gets its own lesson later.
\end{cousinweb}

\subsection*{Your turn}
Say these aloud:

\begin{itemize}
\raggedright
  \item \textquotedblleft{}medius diēs\textquotedblright{} — mid-day, then \textquotedblleft{}media nox\textquotedblright{} — mid-night
  \item the three fates — \textquotedblleft{}Spanish kept both halves, French kept only night, English used a different phrase entirely\textquotedblright{}
  \item what English's real ancestor for \textquotedblleft{}noon\textquotedblright{} meant — \textquotedblleft{}the ninth hour\textquotedblright{}
\end{itemize}

\begin{tcolorbox}[breakable,colback=teal!4,colframe=teal!35!black,title={Before you move on}]
What do \textbf{medius diēs} and \textbf{media nox} literally mean? (\textquotedblleft{}The middle day\textquotedblright{} = noon; \textquotedblleft{}the middle night\textquotedblright{} = midnight.) Which language kept both halves of both phrases fully transparent? (\textbf{Spanish} — \emph{mediodía, medianoche}.) Which language kept \textquotedblleft{}night\textquotedblright{} but not \textquotedblleft{}day\textquotedblright{}? (\textbf{French} — \emph{nuit} stayed alive, \emph{-di} became a fossil.) Does English's \emph{noon} come from this same phrase? (\textbf{No} — from a Latin phrase meaning \textquotedblleft{}the ninth hour,\textquotedblright{} which drifted from afternoon to midday.)
\end{tcolorbox}
